\begin{table*}[t]\centering\caption{Path-MP ablations: MRR (mean $\pm$ std over seeds) and paired difference to the 2-layer model of Table~\ref{tab:main}, in both settings. ``No edge removal'' trains with the query edge present in the message graph.}\label{tab:depth}
\begin{tabular}{lcccc}\toprule
Variant & MRR trans. & MRR induct. & $\Delta$ vs L=2 (trans.) & $\Delta$ vs L=2 (induct.) \\\midrule
1 layer & 0.004$\pm$0.000 & 0.004$\pm$0.000 & -0.984 & -0.985 \\
2 layers (main) & 0.988$\pm$0.004 & 0.989$\pm$0.007 & -- & -- \\
3 layers & 0.989$\pm$0.004 & 0.988$\pm$0.011 & 0.002 & -0.001 \\
2 layers, no edge removal & 0.507$\pm$0.073 & 0.458$\pm$0.090 & -0.480 & -0.531 \\
\bottomrule\end{tabular}

\end{table*}

\begin{table}[t]\centering\caption{Inductive MRR (mean $\pm$ std over seeds) when only a fraction of the new graph's train+validation triples forms the message graph (or fold-in data); the test triples are unchanged. The 100\% column is the main result.}\label{tab:density}
\resizebox{\columnwidth}{!}{\begin{tabular}{lcccc}\toprule
System & 25\% & 50\% & 75\% & 100\% \\\midrule
Path-MP (2-hop) & 0.292$\pm$0.052 & 0.700$\pm$0.041 & 0.934$\pm$0.019 & 0.989$\pm$0.007 \\
TransE+fold-in & 0.329$\pm$0.028 & 0.547$\pm$0.038 & 0.611$\pm$0.036 & 0.644$\pm$0.025 \\
Rule oracle & 0.175$\pm$0.043 & 0.515$\pm$0.048 & 0.803$\pm$0.027 & 0.950$\pm$0.020 \\
\bottomrule\end{tabular}
}
\end{table}

\begin{figure}[t]\centering\includegraphics[width=\columnwidth]{sweep_density_mrr.pdf}
\caption{Inductive MRR against the observed fraction of the new graph (mean $\pm$ std over seeds), from the same runs as Table~\ref{tab:density}.}\label{fig:density}\end{figure}

\vspace{4pt}\noindent\textbf{H4: depth and edge removal (supported).} With one layer the model scores at chance (MRR \rhval{analysis/analysis-final/inductive/lay_trans_l1_mean/mean:3} transductive, \rhval{analysis/analysis-final/inductive/lay_induc_l1_mean/mean:3} inductive; a constant scorer under the even tie-splitting rule gets about $2/n$ with $n$ the number of entities), a paired difference of \rhval{analysis/analysis-final/inductive/lay_induc_l1_vs_l2_diff/mean:3} in the inductive setting. This is expected: with the query edge removed, no answer is within one hop of the query entity for any of the four relations, so the one-layer scores carry no information about the candidates; we did not run a probe that isolates this, so it is the likely, not a demonstrated, reason. A third layer changes nothing measurable: the paired difference to two layers is \rhval{analysis/analysis-final/inductive/lay_trans_l3_vs_l2_diff/mean:3} (paired $p=\rhval{analysis/analysis-final/inductive/lay_trans_l3_vs_l2_p/mean:2}$) transductively and \rhval{analysis/analysis-final/inductive/lay_induc_l3_vs_l2_diff/mean:3} ($p=\rhval{analysis/analysis-final/inductive/lay_induc_l3_vs_l2_p/mean:2}$) inductively; with five seeds this means ``no difference detected'', not equivalence, and the model is at the ceiling so a deeper model has little room to help. Training without removing the query edge is the most damaging change after removing a layer: MRR falls to \rhval{analysis/analysis-final/inductive/nodrop_trans_mean/mean:3} (transductive) and \rhval{analysis/analysis-final/inductive/nodrop_induc_mean/mean:3} (inductive), differences of \rhval{analysis/analysis-final/inductive/nodrop_trans_vs_l2_diff/mean:3} and \rhval{analysis/analysis-final/inductive/nodrop_induc_vs_l2_diff/mean:3} to the main model (paired $p$ below \rhval{const/p_thresh:3}). With the edge present at training time the model can score the answer from the edge itself; at test time that edge is absent by construction. This is the likely mechanism; we did not measure which input the model relies on. The variance across seeds is high for this variant (std \rhval{analysis/analysis-final/inductive/nodrop_induc_std/mean:3} inductive). In the no-removal variant sibling is hurt most (MRR \rhval{analysis/analysis-final/inductive/nodrop_induc_sibling_mean/mean:3} inductive), whereas uncle is hurt least (\rhval{analysis/analysis-final/inductive/nodrop_induc_uncle_mean/mean:3}); we have no isolating run for this ordering.

\vspace{4pt}\noindent\textbf{H5: sparser new graphs (supported for Path-MP; the fold-in baseline degrades too).} Table~\ref{tab:density} and Fig.~\ref{fig:density} vary the observed fraction of the new graph. Path-MP's inductive MRR falls monotonically at the four grid points, from \rhval{analysis/analysis-final/inductive/dens_pm_f100_mean/mean:3} with everything observed to \rhval{analysis/analysis-final/inductive/dens_pm_f75_mean/mean:3}, \rhval{analysis/analysis-final/inductive/dens_pm_f50_mean/mean:3} and \rhval{analysis/analysis-final/inductive/dens_pm_f25_mean/mean:3} at the three sparser levels; the step from the full to the 75\% graph is itself significant (paired difference \rhval{analysis/analysis-final/inductive/dens_pm_f100_vs_f75_diff/mean:3}, $p=\rhval{analysis/analysis-final/inductive/dens_pm_f100_vs_f75_p/mean:sci2}$). The rule oracle falls similarly, so the loss is largely a loss of information about the rules' bodies in the observed graph rather than a defect of the model: a path model cannot answer a query whose supporting facts are not observed. Fold-in TransE degrades more gently (\rhval{analysis/analysis-final/inductive/dens_tf_f100_mean/mean:3} to \rhval{analysis/analysis-final/inductive/dens_tf_f25_mean/mean:3}) and the ranking of the systems changes with density. Path-MP is ahead of fold-in at the 75\% and 50\% levels (paired differences \rhval{analysis/analysis-final/inductive/dens_pm_vs_tf_f75_diff/mean:3} and \rhval{analysis/analysis-final/inductive/dens_pm_vs_tf_f50_diff/mean:3}), but at 25\% fold-in TransE has the higher mean (difference \rhval{analysis/analysis-final/inductive/dens_pm_vs_tf_f25_diff/mean:3} for Path-MP, paired $p=\rhval{analysis/analysis-final/inductive/dens_pm_vs_tf_f25_p/mean:2}$), which with five seeds is inconclusive. Path-MP stays above the oracle at all three sparser levels (paired $p$ below \rhval{const/p_thresh:3}). Note that in this sweep Path-MP is applied without retraining on sparse graphs; training on sparser graphs might narrow the gap, which we did not test.
